%! Author = sriadityavedantam
%! Date = 3/25/26

\section{Homework 2}

\begin{problem}[Subsets of Uncountable Sets]{
    \begin{enumerate}[label={(\alph*)}]
\item Let $C\subseteq[0,1]$ be uncountable.
Show there exists $a\in(0,1)$ such that $C\cap[a,1]$ is uncountable.
\item Let $A$ be the set of all $a\in(0,1)$ such that $C\cap[a,1]$ is uncountable, and define $\alpha=\sup(A)$.
Is $C\cap[\alpha,1]$ uncountable?
\item Does the statement in \textbf{(a)} remain true if “uncountable” is replaced by “infinite”?
    \end{enumerate}
}
    \textbf{(a).}
    For each $n\in\n$, define
    \[
        A_n\coloneqq C\cap\left[\frac{1}{n+1},1\right].
    \]
    Then
    \[
        \bigcup_{n=1}^\infty A_n=C\setminus\bigl(C\cap\{0\}\bigr).
    \]
    The set $C\cap\{0\}$ has at most one element, so if every $A_n$ were countable, then the countable union
    \[
        \bigcup_{n=1}^\infty A_n
    \]
    would be countable, and adjoining at most one point would still give a countable set.
    That would imply $C$ is countable, a contradiction.
    Hence some $A_n$ is uncountable.
    Set
    \[
        a\coloneqq\frac{1}{n+1}\in(0,1).
    \]
    Then
    \[
        C\cap[a,1]
    \]
    is uncountable.

    \textbf{(b).}
    Not necessarily.
    Take
    \[
        C=[0,1).
    \]
    Then for every $a\in(0,1)$,
    \[
        C\cap[a,1]=[a,1)
    \]
    is uncountable, so
    \[
        A=(0,1)
    \]
    and hence
    \[
        \alpha=\sup(A)=1.
    \]
    But then
    \[
        C\cap[\alpha,1]=C\cap[1,1]=\varnothing,
    \]
    which is not uncountable.
    So the answer is no.

    \textbf{(c).}
    No.
    Take
    \[
        C=\left\{\frac{1}{n}:n\in\n\right\}\subset[0,1].
    \]
    Then $C$ is infinite, but for every $a\in(0,1)$,
    \[
        C\cap[a,1]
    \]
    is finite.
    So the statement fails if “uncountable” is replaced by “infinite.”
\end{problem}

\begin{problem}[Arithmetic of Suprema]{
    Let $A,B\subset\rr$ be nonempty sets that are bounded above.
    \begin{enumerate}[label={(\alph*)}]
\item Define
\[
    C\coloneqq\{x:x=a+b\text{ for some }a\in A,\ b\in B\}.
\]
Show that
\[
    \sup(C)=\sup(A)+\sup(B).
\]
\item Show there exists a monotonically increasing sequence $(x_n)$, with all $x_n\in A$, such that $x_n\to\sup(A)$.
    \end{enumerate}
}
    \textbf{(a).}
    Let
    \[
        \alpha\coloneqq\sup(A)
        \qquad\text{and}\qquad
        \beta\coloneqq\sup(B).
    \]
    For all $a\in A$ and $b\in B$,
    \[
        a\leq\alpha
        \qquad\text{and}\qquad
        b\leq\beta.
    \]
    Hence
    \[
        a+b\leq\alpha+\beta.
    \]
    So $\alpha+\beta$ is an upper bound for $C$, and therefore
    \[
        \sup(C)\leq\alpha+\beta.
    \]
    Now let $\eps>0$.
    Since $\alpha=\sup(A)$, there exists $\gamma\in A$ such that
    \[
        \gamma>\alpha-\frac{\eps}{2}.
    \]
    Since $\beta=\sup(B)$, there exists $\delta\in B$ such that
    \[
        \delta>\beta-\frac{\eps}{2}.
    \]
    Then
    \[
        \gamma+\delta>\alpha+\beta-\eps.
    \]
    Because $\gamma+\delta\in C$, we have
    \[
        \sup(C)\geq\gamma+\delta>\alpha+\beta-\eps.
    \]
    Since this holds for all $\eps>0$,
    \[
        \sup(C)\geq\alpha+\beta.
    \]
    Thus
    \[
        \sup(C)=\alpha+\beta.
    \]

    \textbf{(b).}
    Let
    \[
        \alpha\coloneqq\sup(A).
    \]
    For each $n\in\n$, since $\alpha-\frac1n$ is not an upper bound of $A$, there exists $a_n\in A$ such that
    \[
        \alpha-\frac1n<a_n\leq\alpha.
    \]
    Now define
    \[
        x_1\coloneqq a_1,
        \qquad
        x_n\coloneqq\max\{x_{n-1},a_n\}\quad\text{for }n\geq2.
    \]
    Then each $x_n\in A$, and $(x_n)$ is monotonically increasing.
    Also
    \[
        x_n\leq\alpha
        \qquad\text{for all }n.
    \]
    Let $\eps>0$.
    Choose $N\in\n$ such that
    \[
        \frac1N<\eps.
    \]
    Then
    \[
        a_N>\alpha-\eps.
    \]
    Since $x_n\geq a_N$ for all $n\geq N$,
    \[
        \alpha-\eps<x_n\leq\alpha
        \qquad\text{for all }n\geq N.
    \]
    Thus
    \[
        \abs{x_n-\alpha}<\eps
        \qquad\text{for all }n\geq N.
    \]
    Hence
    \[
        x_n\to\alpha=\sup(A).
    \]
\end{problem}

\begin{problem}[Limits of a Root]{
    Let $x_n\geq0$ for all $n\in\n$.
    \begin{enumerate}[label={(\alph*)}]
\item If $(x_n)\to0$, show that $(\sqrt{x_n})\to0$.
\item If $(x_n)\to x$, show that $(\sqrt{x_n})\to\sqrt{x}$.
    \end{enumerate}
}
    \textbf{(a).}
    Let $\eps>0$.
    Because $x_n\to0$, there exists $N\in\n$ such that for all $n\geq N$,
    \[
        \abs{x_n}<\eps^2.
    \]
    Since $x_n\geq0$,
    \[
        0\leq x_n<\eps^2.
    \]
    Taking square roots gives
    \[
        0\leq \sqrt{x_n}<\eps.
    \]
    Hence
    \[
        \abs{\sqrt{x_n}-0}<\eps
    \]
    for all $n\geq N$.
    So
    \[
        \sqrt{x_n}\to0.
    \]

    \textbf{(b).}
    Since $x_n\geq0$ for all $n$ and $x_n\to x$, we must have $x\geq0$.
    Now
    \[
        \abs{\sqrt{x_n}-\sqrt{x}}
        =
        \frac{\abs{x_n-x}}{\sqrt{x_n}+\sqrt{x}}.
    \]
    If $x=0$, then this is part \textbf{(a)}.
    So assume $x>0$.
    Choose $N_1$ such that for $n\geq N_1$,
    \[
        \abs{x_n-x}<\frac{x}{2}.
    \]
    Then for $n\geq N_1$,
    \[
        x_n>\frac{x}{2},
    \]
    so
    \[
        \sqrt{x_n}+\sqrt{x}\geq \sqrt{x}.
    \]
    Hence
    \[
        \abs{\sqrt{x_n}-\sqrt{x}}
        \leq \frac{\abs{x_n-x}}{\sqrt{x}}.
    \]
    Now let $\eps>0$.
    Choose $N_2$ such that for $n\geq N_2$,
    \[
        \abs{x_n-x}<\eps\sqrt{x}.
    \]
    Then for $n\geq \max\{N_1,N_2\}$,
    \[
        \abs{\sqrt{x_n}-\sqrt{x}}<\eps.
    \]
    Therefore
    \[
        \sqrt{x_n}\to\sqrt{x}.
    \]
\end{problem}

\begin{problem}[Examples of Limit Convergence]{
    Using only the definition of convergence, prove that if $x_n\to2$, then
    \begin{enumerate}[label={(\alph*)}]
\item \(\left(\frac{2x_n-1}{3}\right)\to1\),
\item \(\left(\frac{1}{x_n}\right)\to\frac12\).
    \end{enumerate}
}
    \textbf{(a).}
    Let $\eps>0$.
    Since $x_n\to2$, there exists $N\in\n$ such that for all $n\geq N$,
    \[
        \abs{x_n-2}<\frac{3\eps}{2}.
    \]
    Then
    \begin{align*}
        \abs{\frac{2x_n-1}{3}-1}
        &=\abs{\frac{2x_n-4}{3}}\\
        &=\frac23\abs{x_n-2}\\
        &<\frac23\cdot\frac{3\eps}{2}\\
        &=\eps.
    \end{align*}
    Hence
    \[
        \frac{2x_n-1}{3}\to1.
    \]

    \textbf{(b).}
    Let $\eps>0$.
    Since $x_n\to2$, there exists $N_1$ such that for all $n\geq N_1$,
    \[
        \abs{x_n-2}<1.
    \]
    Then for $n\geq N_1$,
    \[
        1<x_n<3,
    \]
    so in particular
    \[
        x_n\geq1.
    \]
    Also choose $N_2$ such that for all $n\geq N_2$,
    \[
        \abs{x_n-2}<2\eps.
    \]
    Now for $n\geq\max\{N_1,N_2\}$,
    \begin{align*}
        \abs{\frac1{x_n}-\frac12}
        &=\abs{\frac{2-x_n}{2x_n}}\\
        &=\frac{\abs{x_n-2}}{2x_n}\\
        &\leq \frac{\abs{x_n-2}}{2}\\
        &<\frac{2\eps}{2}\\
        &=\eps.
    \end{align*}
    Hence
    \[
        \frac1{x_n}\to\frac12.
    \]
\end{problem}

\begin{problem}[Cesàro Means]{
    \begin{enumerate}[label={(\alph*)}]
\item Show that if $(x_n)$ converges to $L$, then the averages
\[
    y_n=\frac{x_1+x_2+\cdots+x_n}{n}
\]
also converge to $L$.
\item Give an example where $(y_n)$ converges even though $(x_n)$ does not.
    \end{enumerate}
}
    \textbf{(a).}
    Assume $x_n\to L$.
    Let
    \[
        y_n=\frac{x_1+\cdots+x_n}{n}.
    \]
    Then
    \[
        y_n-L=\frac{(x_1-L)+\cdots+(x_n-L)}{n}.
    \]
    Let $\eps>0$.
    Choose $N$ such that for all $n\geq N$,
    \[
        \abs{x_n-L}<\eps.
    \]
    Then
    \begin{align*}
        \abs{y_n-L}
        &\leq \frac1n\sum_{k=1}^n \abs{x_k-L}\\
        &= \frac1n\sum_{k=1}^{N-1}\abs{x_k-L}
        +\frac1n\sum_{k=N}^{n}\abs{x_k-L}\\
        &< \frac{C}{n}+\frac{n-N+1}{n}\eps,
    \end{align*}
    where
    \[
        C\coloneqq \sum_{k=1}^{N-1}\abs{x_k-L}.
    \]
    Since $\frac{C}{n}\to0$ and $\frac{n-N+1}{n}\leq1$, it follows that
    \[
        \abs{y_n-L}\to0.
    \]
    Hence
    \[
        y_n\to L.
    \]

    \textbf{(b).}
    Take
    \[
        x_n=(-1)^n.
    \]
    Then $(x_n)$ does not converge.
    But the averages satisfy
    \[
        y_{2k}=0
        \qquad\text{and}\qquad
        y_{2k+1}=-\frac{1}{2k+1}.
    \]
    So
    \[
        y_n\to0.
    \]
\end{problem}